\chapter{Desarrollo y ejecución íntegros en Docker}
\section{Decisión y alcance}
El alumno ha establecido que todo el proyecto se desarrolle y ejecute en contenedores Docker. Se utilizarán contenedores Linux y Docker Compose. Las herramientas locales detectadas anteriormente son inventario, no requisitos para compilar o ejecutar BatchFlow. El host aporta Docker Engine accesible, Compose, Git y el editor; el navegador sigue siendo el cliente del frontend.

También se prevén contenedores de herramientas para Maven/JDK, compilador GnuCOBOL, Node/Angular y LaTeX. No significa arrancar toda la infraestructura para cada test unitario: los perfiles y etapas separan herramientas de servicios de ejecución. Dockerfiles y Compose forman parte del aprendizaje y aún no están creados.

\section{Imágenes y etapas de construcción}
\begin{longtable}{L{0.30\textwidth}L{0.62\textwidth}}
\textbf{Componente} & \textbf{Empaquetado previsto}\\\hline
Servicios Java y gateway & Etapa Maven/JDK para compilar y probar; etapa final con runtime Java y artefactos necesarios.\\
Servicios con COBOL & Compilación nativa en etapa compatible; ejecutable y bibliotecas GnuCOBOL en el runtime del mismo servicio Java.\\
Angular & Node para dependencias, tests y build; servidor de desarrollo en su perfil y servidor HTTP de estáticos para la imagen de entrega.\\
PostgreSQL, Kafka e identidad & Imágenes versionadas, configuración explícita y datos persistentes donde corresponda.\\
Documentación & Imagen LaTeX versionada, fuentes montadas y PDF como salida; auxiliares en área temporal.\\
\end{longtable}

Los builds multietapa permiten excluir compiladores y cachés innecesarios de las imágenes de ejecución. Se comprobará la compatibilidad de bibliotecas nativas: copiar un ejecutable COBOL no garantiza que funcione en una distribución distinta. El contenedor del servicio ejecuta directamente el binario Linux; no invoca wsl.exe. Un proceso hijo COBOL forma parte de ese servicio, no es otro microservicio por usar un lenguaje diferente.

Las versiones concretas y las familias de imágenes se elegirán y registrarán antes de implementar. No se utilizará una etiqueta flotante latest como definición de un entorno reproducible.

\section{Compose, redes y disponibilidad}
La composición tendrá perfiles o configuraciones separadas para desarrollo, pruebas y ejecución de la aplicación. Los contenedores se localizarán por nombres de servicio y puertos internos, no mediante localhost del host. Los puertos publicados se limitarán a los necesarios para navegador, identidad y diagnóstico local documentado.

El navegador no resuelve nombres privados de Compose. La URL pública de la API será accesible desde el host, con proxy de desarrollo o enrutamiento del frontend/gateway para evitar confundir direcciones internas y externas. Kafka e identidad tendrán endpoints coherentes con cada cliente; advertised listeners e issuer OIDC se validarán en una prueba de conexión.

Se definirán healthchecks y dependencias por disponibilidad. depends\_on por sí solo ordena el arranque, pero no garantiza que la base de datos o Kafka estén listos. Las aplicaciones también deben recuperarse de una dependencia que se caiga después del arranque; no se sustituye resiliencia por sleeps fijos.

\section{Persistencia y configuración}
Los datos duraderos de PostgreSQL, Kafka y el proveedor de identidad utilizarán volúmenes nombrados. Las cuentas de base de datos siguen separadas por servicio. Los archivos temporales Java--COBOL se aíslan por ejecución y se limpian según una política explícita; no se confunden con el almacenamiento duradero de un lote.

Configuración no sensible y ejemplos se versionan; secretos reales, archivos locales de entorno y credenciales no se incluyen en Git ni en capas de imagen. Se prevé Compose secrets con permisos por servicio. Debe documentarse que su archivo fuente local también necesita protección: no constituye por sí solo una bóveda de secretos cifrada.

Las imágenes de aplicación ejecutarán con usuario sin privilegios cuando sea posible y escribirán únicamente en rutas previstas. Los logs irán a stdout/stderr, con correlación; la respuesta estructurada de COBOL conservará su canal definido.

\section{Desarrollo, tests y acceso al daemon}
El alumno editará fuentes desde el host y ejecutará compiladores y tests en contenedores. Cachés de Maven y dependencias de Node se ubicarán en volúmenes de herramientas, sin incorporar node\_modules ni target al repositorio. Se controlarán permisos de los archivos generados y se podrán exponer puertos de depuración solo en desarrollo.

Los unit tests no requieren PostgreSQL ni Kafka. Para integración se usarán dependencias reales, con Testcontainers cuando proceda. Ejecutar Testcontainers desde un runner en contenedor exige acceso explícito a un daemon Docker y conectividad hacia los contenedores creados. Esa configuración se limitará al runner de pruebas; ningún servicio de aplicación necesita acceso al socket Docker. Se validará antes de afirmar que los tests de integración están listos.

Se puede disponer de una composición de integración con dependencias explícitas para pruebas que no usen Testcontainers. No se presentará un mock de Kafka o de PostgreSQL como evidencia equivalente a esa integración.

\section{Pruebas de aceptación de la contenerización}
\begin{enumerate}
\item Desde un clon limpio, ejecutar el test de importes en el contenedor de herramientas sin JDK/Maven locales.
\item Compilar COBOL y ejecutar una prueba real Java--COBOL dentro de su runtime Linux.
\item Levantar la composición con healthchecks; completar una operación desde el navegador con URLs correctas.
\item Reiniciar servicios sin perder datos duraderos y observar recuperación ante una dependencia temporalmente no disponible.
\item Construir la imagen de Angular y generar el maestro PDF mediante contenedores de herramientas.
\item Confirmar que Git no incluye secretos, cachés ni archivos de ejecución generados.
\end{enumerate}

Estos son criterios pendientes de implementación, no resultados ya obtenidos. El primer hito será el runner Maven; las imágenes restantes se incorporan cuando exista su incremento funcional. El desarrollo de Dockerfiles y Compose lo realiza el alumno, con revisión y pruebas del tutor.

\section{Referencias de Docker}
\begin{itemize}
\item \href{https://docs.docker.com/build/building/multi-stage/}{Docker: builds multietapa}.
\item \href{https://docs.docker.com/compose/how-tos/startup-order/}{Compose: disponibilidad y orden de arranque}.
\item \href{https://docs.docker.com/compose/how-tos/use-secrets/}{Compose: secretos por servicio}.
\end{itemize}
